\subsubsection{Péndulo simple vía formulación Lagrangiana}

Una partícula de masa $m$ está suspendida de un punto fijo $O$ mediante una cuerda ideal (inextensible, de masa despreciable) de longitud $\ell$, formando un péndulo simple que oscila en un plano vertical, como se muestra en la Figura \ref{fig:c5_e01_p01_pendulo_simple}. Sea $\theta$ el ángulo que forma la cuerda con la vertical.

\begin{enumerate}
    \item Identifique el número de grados de libertad y la coordenada generalizada adecuada.
    \item Construya el Lagrangiano $L = T - V$ del sistema.
    \item Obtenga la ecuación de movimiento mediante la ecuación de Euler-Lagrange y compárela con el resultado newtoniano conocido.
\end{enumerate}

\begin{figure}[!ht]
    \centering
    \begin{tikzpicture}[>=Latex, scale=1.8]
        % --- Parámetros ---
        \def\ell{2.4}     % Longitud de la cuerda
        \def\Theta{32}    % Ángulo respecto a la vertical

        % --- Soporte fijo ---
        \coordinate (O) at (0,0);
        \draw[pattern=north east lines, pattern color=gray!60] (-0.3,0.15) rectangle (0.3,0.35);
        \draw[thick] (-0.35,0.15) -- (0.35,0.15);
        \filldraw (O) circle (1.5pt) node[above] {$O$};

        % --- Vertical de referencia ---
        \draw[dashed, gray] (O) -- (0,-\ell-0.4);

        % --- Cuerda y masa ---
        \coordinate (P) at ({\ell*sin(\Theta)}, {-\ell*cos(\Theta)});
        \draw[thick] (O) -- (P);
        \filldraw[fill=gray!15, thick] (P) circle (3.5pt);
        \node[right] at (P) {$m$};

        % --- Angulo theta ---
        \draw[->] (0,-0.7) arc (-90:{-90+\Theta}:0.7);
        \node at (-90+\Theta/2:0.95) {$\theta$};

        % --- Vector gravedad ---
        \draw[->, thick, myred] (P) -- ++(0,-0.9) node[below, color=myred] {$m\vec{g}$};

        % --- Etiqueta longitud ---
        \node[left] at ({\ell/2*sin(\Theta)}, {-\ell/2*cos(\Theta)}) {$\ell$};

    \end{tikzpicture}
    \caption{Péndulo simple de longitud $\ell$, con coordenada generalizada $\theta$.}
    \label{fig:c5_e01_p01_pendulo_simple}
\end{figure}
